\scanpage{24}
% archive sheet number: 23
\mspage{- 10 -}\hfill n$^{\mathrm{o}}$ 276

en vertu du corollaire précédent on peut donc se ramener\handwritten{/}au cas où $E$ est \uline{libre}, i.e. est produit direct d'un nombre fini de facteurs isomorphes à $A$. Raisonnant par récurrence, on est alors ramené à établir que si deux modules $E$ et $F$ (munis des filtrations associées à la filtration de $A$) sont complets il en est de même du module $E\times F$ : or les relations
\[
A_p E\times A_q F\handwritten{\subset} A_n(E\times F)\handwritten{\subset} A_n E\times A_n F\qquad(p,q\handwritten{\leqslant} n)
\]
montrent que $E\times F$ est le produit direct des groupes \uline{topologiques} $E$ et $F$. il est donc complet si $E$ et $F$ le sont.

PROPOSITION 5 - \uline{Soient} $A$\uline{un anneau commutatif}, $\mathrm{m}$ \uline{un idéal de} $A$, \uline{et supposons l'anneau} $A$ \uline{séparé} (resp. complet) \uline{pour la topologie} $\mathrm{m}$-\uline{adique}. \uline{Alors L'anneau de séries formelles} $A\handwritten{[[}X\handwritten{]]}$, \uline{filtré par les puissances de l'idéal} \inserted{\handwritten{engendré}} \uline{par} $\mathrm{m}$ \uline{et} $X$, \uline{est séparé} (resp. \uline{complet}).
% slanted
\margin{\handwritten{Dire \unsure{que la|que} \unsure{top. en quest.|top. en question|topologie} celle de la conv. simple des coefficients.}}

Soit $\mathrm{n}$ l'idéal de $B={}'A\handwritten{[[}X\handwritten{]]}$ engendré par $\mathrm{m}$ et $X$ : c'est l'ensemble des séries formelles $\handwritten{\sum} a_n X^n$ telles que $a_0\handwritten{\in}\mathrm{m}$. On déduit de là que pour tout entier $p$, l'idéal $\mathrm{n}^p$ est formé des séries $\handwritten{\sum} a_n X^n$ telles que l'on ait $a_n\handwritten{\in}\mathrm{m}^{p-n}$ pour $0\handwritten{\leqslant} n\handwritten{\leqslant} p$. Il résulte immédiatement de là que $B$ est séparé si $A$ est séparé. Supposons de plus $A$ complet, et considérons dans $B$ une série $\handwritten{\sum} y_i$ dont le terme général
\[
y_i=\handwritten{\sum} a_{in}X^n
\]
tend vers 0 ; pour tout entier $p$ on a donc
\begingroup\makeatletter\tagsleft@true\makeatother
\[
a_{in}\handwritten{\in}\mathrm{m}^{p-n}\qquad(0\handwritten{\leqslant} n\handwritten{\leqslant} p)\ \text{pour }i>i_p.\tag*{(1)}
\]
\endgroup

En fixant $n$ et en faisant varier $p$ on voit, puisque $A$ est complet, que la série $\handwritten{\sum}_i a_{in}$ \unsure{c\deleted{n}onverge|converge} dans $A$ ; soit $a_n$ l'une des ses sommes (on ne suppose pas $A$ séparé). En vertu de (1) et du fait que les puissances de $\mathrm{m}$ sont fermées dans $A$, on a, quels que soient $n$ et $p\handwritten{\geqslant} n$, la relation
\[
\handwritten{\sum}_{i>i_p}\unsure{a|e}_{in}\handwritten{\in}\mathrm{m}^{p-n}
\]

